\documentclass[11pt]{article}
\usepackage{fontspec}
\usepackage[french]{babel}
\usepackage{amsmath}
\usepackage{amssymb}
\usepackage{amsthm}
\usepackage{mathrsfs}
\usepackage{tikz-cd}
\usepackage{enumitem}

\newcommand{\printedpage}[1]{\marginpar{\footnotesize p.~#1}}
\newcommand{\msp}[1]{\marginpar{\footnotesize [#1]}}
\newcommand{\Fl}{\mathrm{Fl}}
\newcommand{\Cof}{\mathrm{Cof}}
\newcommand{\deq}{\overset{\text{déf}}{=}}
\newcommand{\univ}{\mathrm{univ}}

\begin{document}

\printedpage{24}
% head: A. Grothendieck : Dérivateurs, Ch. XV, Édition des Universités de Montpellier II et Paris VII
% foot: Version 5 novembre 2010

\msp{page 35}
En effet, par le corollaire 2 de la proposition 4.2, l'hypothèse sur $TC$ implique que $TC = \widetilde{TC}$, donc comme $TC \subset C{\cap}W$, l'hypothèse $C{\cap}W \subset \widetilde{TC}$ $(= TC)$ équivaut à $C{\cap}W = TC$.

\textbf{Corollaire 2.} \textit{Supposons $C$ et $W$ stables par facteurs directs. Alors les conditions équivalentes de la proposition} 4.3 \textit{équivalent à}
\[
C\cap W \;=\; \widetilde{TC}\ . \tag*{(4.3.5)}
\]
\textit{Dualement, si $F$, $W$ sont stables par facteurs directs, les conditions équivalentes de la proposition} 4.3 \textit{équivalent à}
\[
F\cap W \;=\; \overline{TF}\ . \tag*{(4.3.5$'$)}
\]

En effet, si $C$ et $W$ sont stables par facteurs directs, il en est de même de $C\cap W$, et dès lors si $C\cap W\subset\widetilde{TC}$, comme $C\cap W$ contient $TC$, on en conclut $\widetilde{C\cap W}=\widetilde{TC}$. Mais comme $C\cap W\longleftrightarrow TC$ et $C\cap W$ est stable par facteurs directs, $C\cap W$ est clos à gauche par corollaire 2 de la proposition 4.2, \textit{i.e.} $C\cap W=\widetilde{C\cap W}$, d'où $C\cap W=\widetilde{TC}$, \quad \textit{q.e.d.}

\textbf{Corollaire 3.} \textit{Supposons les conditions équivalentes} (a) \textit{à} (c) \textit{de} 4.3 \textit{satisfaites, et de plus} $W$ stable par facteurs directs\footnote[12]{En fait, il suffit de supposer que l'on a $W\supset\widetilde{TC}$, $\overline{TF}$, \textit{i.e.} que tout facteur direct d'un élément de $W$ qui est soit dans $C$, soit dans $F$, est encore dans $W$. Alors $(\widetilde{C},\widetilde{TC},\bar{F},\overline{TF})$ définissent une structure de Quillen close, et on sait que cela implique qu'en fait, $W$ lui aussi est stable par facteurs directs.}. \textit{Alors les conditions préliminaires} (4.3.1) \textit{et} (4.3.2) \textit{restent satisfaites quand on remplace $(C,TF)$ et $(TC,F)$ par leurs clôtures $(C',TF')=(\widetilde{C},\overline{TF})$ et $(TC',F')=(\widetilde{TC},\bar{F})$, et les conditions équivalentes de la proposition} 4.3 \textit{sont encore satisfaites pour ces couples, sous les formes fortes}
\[
\left\{
\begin{aligned}
C'\cap W &\;=\; TC' && (=\widetilde{TC'}\,)\\
F'\cap W &\;=\; TF' && (=\overline{TF'}\,)\\
W &\;=\; TF'\circ TC'\ .
\end{aligned}
\right.
\tag*{(4.3.6)}
\]

\msp{page 36}
\textit{Le triple $(W,C',F')$ est un triple de Quillen clos $($correspondant à une “closed model category” de Quillen, du moins si on suppose $\mathcal{M}$ stable par $\varinjlim$ finies et $\varprojlim$ finies, et que $W$ est modérément saturé$)$.}\footnote[13]{Il en résulte le\par
\textbf{Corollaire 4.} \textit{Soient $W$ un localiseur modérément saturé, et $(C,TF)$, $(TC,F)$ deux précouples de Quillen tels que $TC\subset C\cap W$, $TF\subset F\cap W$. Alors les conditions suivantes sont équivalentes.}
\begin{enumerate}[leftmargin=*]
\item[a)] $TC=C\cap W$.
\item[b)] $TF=F\cap W$.
\item[c)] $W=TF\circ TC$.
\item[d)] $(W,C,F)$ \textit{est un triple de Quillen clos.}
\end{enumerate}}

\textsc{Démonstration.} Les relations (4.3.1) impliquent
\[
\widetilde{TC}\;\subset\;\widetilde{C}\qquad\text{et}\qquad\overline{TF}\;\subset\;\bar{F}\ ,
\]
\printedpage{25}
% head: A. Grothendieck : Dérivateurs, Ch. XV, Édition des Universités de Montpellier II et Paris VII
% foot: Version 5 novembre 2010
\textit{i.e.}
\[
TC'\;\subset\;C'\ ,\qquad TF'\;\subset\;F'\ ,
\]
\textit{i.e.} (4.3.1) pour les nouveaux couples de Quillen. La relation 4.3.2, $W\supset TC,TF$, implique
\[
W\;\supset\;\widetilde{TC}\ ,\ \overline{TF}\ ,\qquad \textit{i.e.}\qquad W\;\supset\;TC'\ ,\ TF'\ ,
\]
en vertu du corollaire 2 à la proposition 4.2 et du fait que $W$ est stable par facteurs directs ; donc on a (4.3.2) pour les nouveaux couples. Comme par hypothèse on a
\[
W\;=\;\overline{TF}\circ\widetilde{TC}\ ,
\]
ceci s'écrit
\[
W\;=\;TF'\circ TC'\ ,
\]
donc les nouveaux couples satisfont aux mêmes conditions, qui cette fois, les couples étant clos, prennent la forme forte (4.3.6).

Prouvons que le triple $(W,C',F')$ satisfait aux axiomes de Quillen d'un triple de Quillen clos. On a bien, par (4.3.6),
\[
\left\{
\begin{aligned}
C' &\;=\;(TF')^*\;=\;(F'\cap W)^*\ , &\qquad F' &\;=\;(TC')_*\;=\;(C'\cap W)_*\ ,\\
W &\;=\;TF'\circ TC'\;=\;C'_*\circ F'^*\ ,
\end{aligned}
\right.
\]
qui correspond à la définition des catégories de modèles de Quillen \textit{clos} : les axiomes M\,1, M\,2, M\,3, M\,4 de Quillen sont réalisés de façon triviale (M\,3 et M\,4 parce que
\msp{page 37}
$C$ et $F$, $TC$ et $TF$ sont clos à gauche resp. à droite). L'axiome M\,5 est la saturation modérée de $W$.

\textbf{4.4.} Voici l'idée pour construire plus ou moins “toutes” les structures de Quillen closes $(C,TC,F,TF)$ associées à un localiseur donné $W\subset\mathcal{M}$. On suppose que $\mathcal{M}$ satisfait aux conditions du paragraphe précédent, concernant l'existence de $\varinjlim$ filtrantes, d'une filtration cardinale, l'accessibilité dans $\mathcal{M}$ (celle des éléments de $\mathrm{Filt}^{\pi_0}$). On part d'une partie petite
\[
C_0\;\subset\;\Fl(\mathcal{M})
\]
telle que
\[
C_{0*}\;\subset\;W\ , \tag*{(4.4.1)}
\]
et on pose
\[
\left\{
\begin{aligned}
C &\;=\;\widetilde{C_0}\ , &\qquad TC &\;=\;C\cap W\ ,\\
TF &\;=\;C_*\;=\;C_{0*}\ , &\qquad F &\;=\;(TC)_*\ .
\end{aligned}
\right.
\tag*{(4.4.2)}
\]
Il y a un problème seulement, vu que l'on sait déjà que $(C,TF)$ est un couple de Quillen par le théorème de factorisation dans §1 (p. 10) : il faut la factorisation pour le couple $(TC,F\;=\;(TC)_*)$. À ce moment là, les conditions préliminaires (4.3.1), (4.3.2) de la proposition 4.3 sont satisfaites :
\[
TC\;\subset C\ ,\qquad TF\;\subset F\ ,\qquad W\;\supset\;TC\ ,\ TF\ ,
\]
ainsi que la condition $(a)$, $C\cap W\subset\widetilde{TC}$, puisque $C\cap W=TC$. Donc on trouve, par l'équivalence des conditions $(a)$, $(a')$, $(c)$, et vu que $TC$ et $TF$ sont clos,
\[
TF\;=\;F\cap W\ ,\qquad W\;=\;TF\circ TC\ ,
\]
\printedpage{26}
% head: A. Grothendieck : Dérivateurs, Ch. XV, Édition des Universités de Montpellier II et Paris VII
% foot: Version 5 novembre 2010
\msp{page 38}
d'autre part l'hypothèse sur $W$ implique que $W$ est stable par facteurs directs, $C$ l'est aussi, étant close à gauche, donc $TC=C{\cap}W$ aussi, donc par proposition 4.1, corollaire 2, on a $TC\;=\;\widetilde{TC}$, $TC$ est close à gauche. On gagne donc, par la dernière assertion du corollaire 3, proposition 4.3 (p. 35).

En fait, indépendamment des conditions particulières sur $\mathcal{M}$, $W$, on vient de prouver ceci :

\textbf{Lemme 4.4.1} (corollaire de la proposition 4.3.). \textit{Soient $\mathcal{M}$ une catégorie, $W$ un localiseur dans $\mathcal{M}$ modérément saturé et stable par facteurs directs. Alors les structures de Quillen closes $(C,F)$ associées à $W$ s'obtiennent ainsi : on prend une partie $C_0$ de} $\Fl(\mathcal{M})$ \textit{telle que}
\[
TF\;\deq\;C_{0*}\;\subset\;W\ ,
\]
\textit{on pose}
\[
C\;=\;\widetilde{C_0}\;=\;TF^*\ ,\qquad TC\;=\;C\cap W\ ,\qquad F\;=\;(TC)_*\ ,
\]
\textit{et on doit se débrouiller pour que $(C,TF)$ et $(TC,F)$ satisfassent à la condition de factorisation de Quillen,} i.e. \textit{soient des précouples de Quillen.} (En fait, $C$, $TF$ et $F$ étant clos par construction, et $TC=C\cap W$ étant stable par facteurs directs ($C$ et $W$ l'étant), ce seront même des \textit{couples} de Quillen.)

\msp{page 39}
Les hypothèses supplémentaires faites tantôt sur $\mathcal{M}$ servaient à assurer que si je prends $C_0$ petite, alors la factorisation marche pour $(C,TF)$ automatiquement. Il reste le problème pour $TC$, $F$.

Notons que si elle marche, alors $TC=\widetilde{TC}=F^*$, donc $TC$ est clos à gauche, et \textit{a fortiori} satisfait les conditions de stabilité correspondantes :
\[
\left\{
\begin{tabular}{rp{0.6\textwidth}}
\textit{a}) & Stable par cochangement de base, par composition, par facteurs directs (ces deux dernières sont assurées automatiquement, car $C$, $W$ sont stables).\\[1ex]
\textit{b}) & Stable par $\varinjlim$ ordinales strictes.
\end{tabular}
\right.
\tag*{(4.4.3)}
\]

L'hypothèse de stabilité sur $W$ sert (entre autres) à garantir la stabilité $(b)$ pour $TC$ (en fait, il suffirait pour cela seulement que $W$ soit stable par $\varinjlim$ ordinales strictes). Il reste la stabilité de $TC=C{\cap}W$ par cochangement de base, qui est problématique car $W$ n'est en général \textit{pas} stable par cochangement de base. C'est pour assurer cette stabilité que je me vois pratiquement contraint ici à imposer
\[
C_0\;\subset\;\Cof_W\ , \tag*{(4.4.4)}
\]
je dis qu'avec l'hypothèse de stabilité faite sur $W$, cela implique
\[
C\;\deq\;\widetilde{C_0}\subset\Cof_W\ . \tag*{(4.4.5)}
\]
\printedpage{27}
% head: A. Grothendieck : Dérivateurs, Ch. XV, Édition des Universités de Montpellier II et Paris VII
% foot: Version 5 novembre 2010
\msp{page 40}
En effet, on a :

\textbf{Proposition 4.4.2.} \textit{Soit $\mathcal{M}$ stable par petites $\varinjlim$ inductives filtrantes, et $W\subset$} $\Fl(\mathcal{M})$ \textit{un localiseur $($modérément saturé$)$} stable par lesdites limites. \textit{Alors} $\Cof_W$ \textit{satisfait les conditions de stabilité} (4.4.3) (“\textit{stabilité de Quillen gauche}”), \textit{et est stable aussi par $\varinjlim$ filtrantes quelconques.}

En effet, la condition $(a)$ est satisfaite sans condition supplémentaire sur $W$. D'autre part, il est immédiat que cette condition par rapport aux $\varinjlim$ filtrantes implique que $\Cof_W$ \textit{est stable par limites inductives filtrantes} dans $\underline{\Fl}(\mathcal{M})$.

\textbf{Corollaire 4.4.3.} \textit{Sous les conditions de la proposition} 4.4.2, $\Cof_W\cap W=W^{\univ}$ $($\textit{ensemble des $W$-équivalences co-universelles dans $\mathcal{M}$}$)$ \textit{est stable par limites inductives filtrantes, et donc “stable à gauche” en sens de} (4.4.3).

\textbf{Corollaire 4.4.4.} \textit{Sous les conditions de la proposition} 4.4.2, \textit{supposons de plus que les sommes amalgamées simples existent dans $\mathcal{M}$, et soit $C_0$ une petite partie de} $\Cof_W$ \textit{dont les éléments soient accessibles, et considérons}
\[
C\;=\;\widetilde{C_0}\ ,\qquad TC\;\deq\;C\cap W\ .
\]
\textit{Alors on a}
\[
C\;=\;\Cof_W\ ,\qquad TC\;=\;C\cap W^{\univ}\;\subset\;W^{\univ}\ ,
\]
\msp{page 41}
\textit{et $C$ et $TC$ sont $Q$-stables à gauche.}

En effet, par le théorème 1 du paragraphe 1 (p. 10) et son corollaire, $\widetilde{C_0}$ est la plus petite partie de $\Fl(\mathcal{M})$ contenant $C_0$ et Q-stable à gauche. Comme $\Cof_W$ est Q-stable à gauche et contient $C_0$, on a donc $C=\widetilde{C_0}\subset\Cof_W$. Donc $TC\deq C{\cap}W\subset\Cof_W{\cap}W\deq W^{\univ}$, donc $TC=C\cap W^{\univ}$. Comme $W^{\univ}$ est lui aussi $Q$-stable à gauche (corollaire précédent), il en est de même de $C\cap W$, \quad \textit{q.e.d.}

\textbf{Corollaire 4.4.5.} \textit{Supposons que tout objet de $\mathcal{M}$ soit accessible. Sous les conditions du corollaire précédent, posant}
\[
TF\;=\;C_{0*}\;=\;C_*\ ,\qquad F\;=\;(TC)_*\ ,
\]
\textit{le couple $(C,TF)$ est un couple de Quillen. Pour qu'il en soit de même du couple $(TC,F)$, il suffit qu'il existe une} petite \textit{partie $TC_0$ de $C\cap W$, telle que $\widetilde{TC_0}=TC$ $(\deq C\cap W)$.}

La première assertion est, on l'a vu tantôt, le théorème 1 du paragraphe 1. La deuxième résulte de ce même théorème. On peut préciser dans ce sens :

\textbf{Proposition 4.4.6.} \textit{Soit $\mathcal{M}$ une catégorie stable par petites $\varinjlim$ filtrantes, par sommes amalgamées, et dont tous les objets sont accessibles, $W$ un localiseur}
\msp{page 42}
\textit{modérément saturé dans $\mathcal{M}$, stable par $\varinjlim$ filtrantes, $C_0$ une} petite \textit{partie de} $\Fl(\mathcal{M})$, \textit{telle que}
\[
C_0\;\subset\;\Cof_W\ ,\qquad \underbrace{C_{0*}}_{(=TF)}\;\subset\;W\ ,
\]
\printedpage{28}
% head: A. Grothendieck : Dérivateurs, Ch. XV, Édition des Universités de Montpellier II et Paris VII
% foot: Version 5 novembre 2010
\textit{posons}
\[
C\;=\;\widetilde{C_0}\qquad TC\;=\;C\cap W\ ,\qquad TF\;=\;C_*\ ,\qquad F\;=\;(TC)_*\ ,
\]
\textit{et soit $TC_0$ une petite partie de $TC=C\cap W$. On a alors}
\[
\left\{
\begin{aligned}
C &\;\subset\;\Cof_W\ , &\qquad TC &\;=\;C\cap W^{\univ}\quad(\textit{pour mémoire})\\
\widetilde{TC_0} &\;\subset\;TC\ , & F_0 &\supset F\quad(\textit{ce dernier trivial})\ ,
\end{aligned}
\right.
\tag*{(4.4.6)}
\]
\textit{et les conditions suivantes sont équivalentes :}
\begin{enumerate}[leftmargin=*]
\item[a)] $\widetilde{TC_0}=TC\quad(\textit{on a}\subset)$\,;
\item[a$'$)] $F_0=F\quad(\textit{on a}\supset)$\,;
\item[b)] $TF=F_0\cap W\quad(\textit{on a}\subset)$\footnote[14]{\textbf{N.B.} On voit facilement que $\boxed{TF=F\cap W,}$ mais ça ne suffit pas pour avoir $TF=F_0\cap W$ pour $TC_0$ convenable.}\,;
\item[c)] $W=TF\circ\widetilde{TC_0}\quad(\textit{on a}\supset)$\footnote[15]{\textbf{N.B.} On sait que $\boxed{W=TF\circ TC,}$ ce qui n'est nullement suffisant pour impliquer $(c)$\,!}\,;
\end{enumerate}
\textit{et impliquent que $(W,C,F)$ est un triple de Quillen clos.}

\textsc{Démonstration.} L'inclusion $\widetilde{TC_0}\;\subset\;TC$ résulte de la Q-stabilité à gauche de $TC$ (cor. 4.4.4) et du fait que $TC_0\subset TC$ est petite (on applique alors le théorème 1 tout comme pour $C\subset\Cof_W\ldots$). L'équivalence de $(a)$, $(b)$, $(c)$ n'est autre que la proposition 4.3 appliquée à $W$ et aux deux couples de Quillen $(C,TF)$, $(\widetilde{TC_0},F_0)$ (compte tenu du théorème 1 du paragraphe 1). Qu'on a alors un triple de Quillen clos a été vu dans le corollaire 3 de la proposition 4.3 (p. 35). Il reste à prouver l'équivalence de $(a)$ et de $(a')$. On a $(a)\Longrightarrow(a')$, puisque $(TC_0)_*=(\widetilde{TC_0})_*=F_0$. Inversement, $(a')$ implique $\widetilde{TC_0}=\widetilde{TC}$, donc le théorème de factorisation pour $(TC,F)$ en vertu du théorème de factorisation, corollaire 1 (p. 11) appliqué à $TC_0\subset TC\subset\widetilde{TC_0}$ et à $F_0=F=(TC_0)_*$.

\msp{page 43}
Ainsi, le seul problème, une fois trouvé $C_0$ (petite) satisfaisant $C_0\subset\Cof_W$, $C_{0*}\subset W$, c'est de trouver une \textit{petite} partie $TC_0\subset TC$ qui “engendre $TC\;=\;C\cap W$ ”, \textit{i.e.} qui satisfasse aux conditions équivalentes $(a)$, $(a')$, $(b)$, $(c)$ de la proposition précédente. J'ai passé des jours à essayer de prouver l'existence de $TC_0$, en supposant $\mathcal{M}$ munie d'une filtration cardinale, et $W$ une partie accessible de $\mathcal{M}$, sans y être arrivé. Bien sûr, tout ce qui nous importe vraiment, c'est que $(TC,F)$ satisfasse à la condition de factorisation. Mais je ne vois aucun moyen de le prouver, si ce n'est en “engendrant” à gauche $TC$ par un \textit{petit} sous-ensemble $TC_0$.

\section*{5\quad Résultats techniques préliminaires sur les parties accessibles de $\mathrm{Ob}\,\mathcal{M}$, stables par limites inductives filtrantes}

Je crois pourtant que je peux arriver à quelque chose d'utilisable, en travaillant systématiquement avec des $C$ ($\subset\Cof_W$, Q-closes à gauche, Q-engendrées à gauche par une \textit{petite} partie $C_0$ de $C$) qui sont stables par toutes limites inductives filtrantes.

\end{document}
